
\documentclass{article}
\usepackage{colortbl}
\usepackage{makecell}
\usepackage{multirow}
\usepackage{supertabular}

\begin{document}

\newcounter{utterance}

\centering \large Interaction Transcript for game `hot\_air\_balloon', experiment `air\_balloon\_survival\_de\_negotiation\_hard', episode 5 with deepseek{-}chat{-}v3.1{-}t1.0.
\vspace{24pt}

{ \footnotesize  \setcounter{utterance}{1}
\setlength{\tabcolsep}{0pt}
\begin{supertabular}{c@{$\;$}|p{.15\linewidth}@{}p{.15\linewidth}p{.15\linewidth}p{.15\linewidth}p{.15\linewidth}p{.15\linewidth}}
    \# & $\;$A & \multicolumn{4}{c}{Game Master} & $\;\:$B\\
    \hline

    \theutterance \stepcounter{utterance}  
    & & \multicolumn{4}{p{0.6\linewidth}}{
        \cellcolor[rgb]{0.9,0.9,0.9}{
            \makecell[{{p{\linewidth}}}]{
                \texttt{\tiny{[P1$\langle$GM]}}
                \texttt{Du nimmst an einem kollaborativen Verhandlungspiel Teil.} \\
\\ 
\texttt{Zusammen mit einem anderen Teilnehmer musst du dich auf eine Reihe von Gegenständen entscheiden, die behalten werden. Jeder von euch hat eine persönliche Verteilung über die Wichtigkeit der einzelnen Gegenstände. Jeder von euch hat eine eigene Meinung darüber, wie wichtig jeder einzelne Gegenstand ist (Gegenstandswichtigkeit). Du kennst die Wichtigkeitsverteilung des anderen Spielers nicht. Zusätzlich siehst du, wie viel Aufwand jeder Gegenstand verursacht.  } \\
\texttt{Ihr dürft euch nur auf eine Reihe von Gegenständen einigen, wenn der Gesamtaufwand der ausgewählten Gegenstände den Maximalaufwand nicht überschreitet:} \\
\\ 
\texttt{Maximalaufwand = 3386} \\
\\ 
\texttt{Hier sind die einzelnen Aufwände der Gegenstände:} \\
\\ 
\texttt{Aufwand der Gegenstände = \{"C33": 313, "A89": 720, "A48": 175, "A04": 461, "A41": 635, "C64": 685, "C49": 544, "A61": 203, "C04": 369, "B53": 539, "B85": 4, "B61": 695, "A45": 399, "A03": 594, "B21": 437\}} \\
\\ 
\texttt{Hier ist deine persönliche Verteilung der Wichtigkeit der einzelnen Gegenstände:} \\
\\ 
\texttt{Werte der Gegenstandswichtigkeit = \{"C33": 138, "A89": 583, "A48": 65, "A04": 262, "A41": 121, "C64": 508, "C49": 461, "A61": 484, "C04": 668, "B53": 389, "B85": 215, "B61": 97, "A45": 500, "A03": 30, "B21": 400\}} \\
\\ 
\texttt{Ziel:} \\
\\ 
\texttt{Dein Ziel ist es, eine Reihe von Gegenständen auszuhandeln, die dir möglichst viel bringt (d. h. Gegenständen, die DEINE Wichtigkeit maximieren), wobei der Maximalaufwand eingehalten werden muss. Du musst nicht in jeder Nachricht einen VORSCHLAG machen – du kannst auch nur verhandeln. Alle Taktiken sind erlaubt!} \\
\\ 
\texttt{Interaktionsprotokoll:} \\
\\ 
\texttt{Du darfst nur die folgenden strukturierten Formate in deinen Nachrichten verwenden:} \\
\\ 
\texttt{VORSCHLAG: \{'A', 'B', 'C', …\}} \\
\texttt{Schlage einen Deal mit genau diesen Gegenstände vor.} \\
\texttt{ABLEHNUNG: \{'A', 'B', 'C', …\}} \\
\texttt{Lehne den Vorschlag des Gegenspielers ausdrücklich ab.} \\
\texttt{ARGUMENT: \{'...'\}} \\
\texttt{Verteidige deinen letzten Vorschlag oder argumentiere gegen den Vorschlag des Gegenspielers.} \\
\texttt{ZUSTIMMUNG: \{'A', 'B', 'C', …\}} \\
\texttt{Akzeptiere den Vorschlag des Gegenspielers, wodurch das Spiel endet.} \\
\texttt{STRATEGISCHE ÜBERLEGUNGEN: \{'...'\}} \\
\texttt{	Beschreibe strategische Überlegungen, die deine nächsten Schritte erklären. Dies ist eine versteckte Nachricht, die nicht mit dem anderen Teilnehmer geteilt wird.} \\
\\ 
\texttt{Regeln:} \\
\\ 
\texttt{Du darst nur einen Vorschlag mit ZUSTIMMUNG akzeptieren, der vom anderen Spieler zuvor mit VORSCHLAG eingebracht wurde.} \\
\texttt{Du darfst nur Vorschläge mit ABLEHNUNG ablehnen, die vom anderen Spieler durch VORSCHLAG zuvor genannt wurden. } \\
\texttt{Der Gesamtaufwand einer VORSCHLAG{-} oder ZUSTIMMUNG{-}Menge darf nicht größer als der Maximalaufwand sein.  } \\
\texttt{Offenbare deine versteckte Wichtigkeitsverteilung nicht.} \\
\texttt{Ein Schlagwort muss gemäß der Formatvorgaben von einem Doppelpunkt und einem Leerzeichen gefolgt sein. Das Argument muss eine Python{-}Menge mit 0 oder mehr Strings sein.  } \\
\texttt{Die Form muss also SCHLAGWORT: \{...\} sein.} \\
\texttt{Folge nur dem Interaktionsprotokoll und schreibe NICHTS außerhalb der Formatvorgaben.} \\
\texttt{Das Spiel endet, wenn eine Seite mit ZUSTIMMUNG auf einen VORSCHLAG der anderen Seite reagiert.  } \\
\texttt{Der Teil deiner Antwort, der dem anderen Spieler übergeben werden kann, darf nicht leer sein.  } \\
\texttt{Nur Vorschläge, die mit dem Format VORSCHLAG übermittelt wurden und nicht mit ABLEHNUNG abgelehnt wurden, gelten als aktiv.  } \\
\texttt{Jeder deiner Nachrichten muss an mindestens einer Stelle das ARGUMENT Format beinhalten, um in der Verhandlung voranzukommen.} \\
\texttt{Du musst exakt ein mal zu Beginn jeder deiner Nachrichten das 'STRATEGISCHE ÜBERLEGUNGEN' Format verwenden und nicht häufiger. Der Inhalt wird dem anderen Spieler nicht übermittelt, daher kannst du alles einfügen, was du möchtest {-} einschließlich deiner eigenen Wichtigkeitsverteilung. Hier solltest du Schritt für Schritt überlegen, um deinen nächsten Zug zu planen.} \\
            }
        }
    }
    & \\ \\

    \theutterance \stepcounter{utterance}  
    & \multicolumn{4}{p{0.6\linewidth}}{
        \cellcolor[rgb]{0.9,0.9,0.9}{
            \makecell[{{p{\linewidth}}}]{
                \texttt{\tiny{[P1$\rangle$GM]}}
                \texttt{STRATEGISCHE ÜBERLEGUNGEN: \{'Ich muss meine höchstwertigen Gegenstände identifizieren, die relativ geringen Aufwand verursachen. C04 hat den höchsten Wert (668) mit moderatem Aufwand (369). A89 hat Wert 583 mit hohem Aufwand 720. C64 hat Wert 508 mit Aufwand 685. A45 hat Wert 500 mit Aufwand 399. B21 hat Wert 400 mit Aufwand 437. B53 hat Wert 389 mit Aufwand 539. A61 hat Wert 484 mit niedrigem Aufwand 203. C49 hat Wert 461 mit Aufwand 544. B85 hat niedrigen Wert 215 aber sehr niedrigen Aufwand 4. Ich sollte auch Gegenstände mit niedrigem Wert für mich aber hohem Aufwand opfern, um Raum für wertvollere Gegenstände zu schaffen. A03 hat nur Wert 30 mit Aufwand 594 {-} sehr schlechtes Verhältnis. B61 hat Wert 97 mit hohem Aufwand 695. A48 hat Wert 65 mit Aufwand 175. A41 hat Wert 121 mit Aufwand 635. C33 hat Wert 138 mit Aufwand 313. A04 hat Wert 262 mit Aufwand 461. Ich sollte versuchen, C04, A61, A45, B21, B85 und evtl. C49 oder A89 zu bekommen, während ich niedrigwertige Gegenstände opfere. Der Maximalaufwand ist 3386, also muss ich eine Kombination wählen, die unter diesem Limit bleibt.'\}} \\
\\ 
\texttt{ARGUMENT: \{'Ich schlage einen fairen Deal vor, der sowohl für mich als auch für dich vorteilhaft sein sollte. Mein Vorschlag berücksichtigt verschiedene Faktoren, um eine ausgewogene Lösung zu finden.'\}} \\
\\ 
\texttt{VORSCHLAG: \{'C04', 'A61', 'A45', 'B21', 'B85', 'C49', 'A48'\}} \\
            }
        }
    }
    & & \\ \\

    \theutterance \stepcounter{utterance}  
    & & \multicolumn{4}{p{0.6\linewidth}}{
        \cellcolor[rgb]{0.9,0.9,0.9}{
            \makecell[{{p{\linewidth}}}]{
                \texttt{\tiny{[GM$\rangle$P2]}}
                \texttt{Du nimmst an einem kollaborativen Verhandlungspiel Teil.} \\
\\ 
\texttt{Zusammen mit einem anderen Teilnehmer musst du dich auf eine Reihe von Gegenständen entscheiden, die behalten werden. Jeder von euch hat eine persönliche Verteilung über die Wichtigkeit der einzelnen Gegenstände. Jeder von euch hat eine eigene Meinung darüber, wie wichtig jeder einzelne Gegenstand ist (Gegenstandswichtigkeit). Du kennst die Wichtigkeitsverteilung des anderen Spielers nicht. Zusätzlich siehst du, wie viel Aufwand jeder Gegenstand verursacht.  } \\
\texttt{Ihr dürft euch nur auf eine Reihe von Gegenständen einigen, wenn der Gesamtaufwand der ausgewählten Gegenstände den Maximalaufwand nicht überschreitet:} \\
\\ 
\texttt{Maximalaufwand = 3386} \\
\\ 
\texttt{Hier sind die einzelnen Aufwände der Gegenstände:} \\
\\ 
\texttt{Aufwand der Gegenstände = \{"C33": 313, "A89": 720, "A48": 175, "A04": 461, "A41": 635, "C64": 685, "C49": 544, "A61": 203, "C04": 369, "B53": 539, "B85": 4, "B61": 695, "A45": 399, "A03": 594, "B21": 437\}} \\
\\ 
\texttt{Hier ist deine persönliche Verteilung der Wichtigkeit der einzelnen Gegenstände:} \\
\\ 
\texttt{Werte der Gegenstandswichtigkeit = \{"C04": 30, "A89": 65, "C64": 97, "A45": 121, "A61": 138, "C49": 215, "B21": 262, "B53": 389, "A04": 400, "B85": 461, "C33": 484, "A41": 500, "B61": 508, "A48": 583, "A03": 668\}} \\
\\ 
\texttt{Ziel:} \\
\\ 
\texttt{Dein Ziel ist es, eine Reihe von Gegenständen auszuhandeln, die dir möglichst viel bringt (d. h. Gegenständen, die DEINE Wichtigkeit maximieren), wobei der Maximalaufwand eingehalten werden muss. Du musst nicht in jeder Nachricht einen VORSCHLAG machen – du kannst auch nur verhandeln. Alle Taktiken sind erlaubt!} \\
\\ 
\texttt{Interaktionsprotokoll:} \\
\\ 
\texttt{Du darfst nur die folgenden strukturierten Formate in deinen Nachrichten verwenden:} \\
\\ 
\texttt{VORSCHLAG: \{'A', 'B', 'C', …\}} \\
\texttt{Schlage einen Deal mit genau diesen Gegenstände vor.} \\
\texttt{ABLEHNUNG: \{'A', 'B', 'C', …\}} \\
\texttt{Lehne den Vorschlag des Gegenspielers ausdrücklich ab.} \\
\texttt{ARGUMENT: \{'...'\}} \\
\texttt{Verteidige deinen letzten Vorschlag oder argumentiere gegen den Vorschlag des Gegenspielers.} \\
\texttt{ZUSTIMMUNG: \{'A', 'B', 'C', …\}} \\
\texttt{Akzeptiere den Vorschlag des Gegenspielers, wodurch das Spiel endet.} \\
\texttt{STRATEGISCHE ÜBERLEGUNGEN: \{'...'\}} \\
\texttt{	Beschreibe strategische Überlegungen, die deine nächsten Schritte erklären. Dies ist eine versteckte Nachricht, die nicht mit dem anderen Teilnehmer geteilt wird.} \\
\\ 
\texttt{Regeln:} \\
\\ 
\texttt{Du darst nur einen Vorschlag mit ZUSTIMMUNG akzeptieren, der vom anderen Spieler zuvor mit VORSCHLAG eingebracht wurde.} \\
\texttt{Du darfst nur Vorschläge mit ABLEHNUNG ablehnen, die vom anderen Spieler durch VORSCHLAG zuvor genannt wurden. } \\
\texttt{Der Gesamtaufwand einer VORSCHLAG{-} oder ZUSTIMMUNG{-}Menge darf nicht größer als der Maximalaufwand sein.  } \\
\texttt{Offenbare deine versteckte Wichtigkeitsverteilung nicht.} \\
\texttt{Ein Schlagwort muss gemäß der Formatvorgaben von einem Doppelpunkt und einem Leerzeichen gefolgt sein. Das Argument muss eine Python{-}Menge mit 0 oder mehr Strings sein.  } \\
\texttt{Die Form muss also SCHLAGWORT: \{...\} sein.} \\
\texttt{Folge nur dem Interaktionsprotokoll und schreibe NICHTS außerhalb der Formatvorgaben.} \\
\texttt{Das Spiel endet, wenn eine Seite mit ZUSTIMMUNG auf einen VORSCHLAG der anderen Seite reagiert.  } \\
\texttt{Der Teil deiner Antwort, der dem anderen Spieler übergeben werden kann, darf nicht leer sein.  } \\
\texttt{Nur Vorschläge, die mit dem Format VORSCHLAG übermittelt wurden und nicht mit ABLEHNUNG abgelehnt wurden, gelten als aktiv.  } \\
\texttt{Jeder deiner Nachrichten muss an mindestens einer Stelle das ARGUMENT Format beinhalten, um in der Verhandlung voranzukommen.} \\
\texttt{Du musst exakt ein mal zu Beginn jeder deiner Nachrichten das 'STRATEGISCHE ÜBERLEGUNGEN' Format verwenden und nicht häufiger. Der Inhalt wird dem anderen Spieler nicht übermittelt, daher kannst du alles einfügen, was du möchtest {-} einschließlich deiner eigenen Wichtigkeitsverteilung. Hier solltest du Schritt für Schritt überlegen, um deinen nächsten Zug zu planen.} \\
\\ 
\texttt{Du erhältst nun die erste Nachricht des anderen Spielers.} \\
\\ 
\texttt{ARGUMENT: \{'Ich schlage einen fairen Deal vor, der sowohl für mich als auch für dich vorteilhaft sein sollte. Mein Vorschlag berücksichtigt verschiedene Faktoren, um eine ausgewogene Lösung zu finden.'\}} \\
\\ 
\texttt{VORSCHLAG: \{'C04', 'A61', 'A45', 'B21', 'B85', 'C49', 'A48'\}} \\
            }
        }
    }
    & \\ \\

    \theutterance \stepcounter{utterance}  
    & & & \multicolumn{4}{p{0.6\linewidth}}{
        \cellcolor[rgb]{0.9,0.9,0.9}{
            \makecell[{{p{\linewidth}}}]{
                \texttt{\tiny{[GM$\langle$P2]}}
                \texttt{STRATEGISCHE ÜBERLEGUNGEN: \{'Der Gegner hat einen Vorschlag mit den Gegenständen C04, A61, A45, B21, B85, C49, A48 gemacht. Ich berechne den Gesamtaufwand: 369 + 203 + 399 + 437 + 4 + 544 + 175 = 2131, was unter dem Maximalaufwand von 3386 liegt. Mein eigener Wert für diesen Vorschlag: C04=30, A61=138, A45=121, B21=262, B85=461, C49=215, A48=583. Summe = 30+138+121+262+461+215+583 = 1810. Ich sollte meine Wichtigkeit maximieren, die maximal erreichbare Summe ist 668+583+508+500+484+461+400+389+262+215+138+121+97+65+30 = 4912, aber das überschreitet den Aufwand. Ich muss einen besseren Deal für mich aushandeln, der mehr meiner hochwertigen Gegenstände enthält, wie A03 (668), B61 (508), A41 (500), C33 (484), A04 (400), B53 (389). Ich lehne diesen Vorschlag ab, da er mir nur 1810 bringt, und mache einen Gegenvorschlag mit höherem Wert für mich.'\}} \\
\\ 
\texttt{ARGUMENT: \{'Dein Vorschlag enthält einige wertvolle Gegenstände, aber ich habe andere Prioritäten. Viele der für mich wichtigen Items sind nicht enthalten. Lass uns einen Deal finden, der für beide Seiten vorteilhafter ist.'\}} \\
\\ 
\texttt{ABLEHNUNG: \{'C04', 'A61', 'A45', 'B21', 'B85', 'C49', 'A48'\}} \\
\\ 
\texttt{VORSCHLAG: \{'A03', 'B61', 'A41', 'C33', 'A04', 'B53', 'B85', 'A61'\}} \\
            }
        }
    }
     \\ \\

    \theutterance \stepcounter{utterance}  
    & & \multicolumn{4}{p{0.6\linewidth}}{
        \cellcolor[rgb]{0.9,0.9,0.9}{
            \makecell[{{p{\linewidth}}}]{
                \texttt{\tiny{[P1$\langle$GM]}}
                \texttt{ARGUMENT: \{'Dein Vorschlag enthält einige wertvolle Gegenstände, aber ich habe andere Prioritäten. Viele der für mich wichtigen Items sind nicht enthalten. Lass uns einen Deal finden, der für beide Seiten vorteilhafter ist.'\}} \\
\\ 
\texttt{ABLEHNUNG: \{'C04', 'A61', 'A45', 'B21', 'B85', 'C49', 'A48'\}} \\
\\ 
\texttt{VORSCHLAG: \{'A03', 'B61', 'A41', 'C33', 'A04', 'B53', 'B85', 'A61'\}} \\
            }
        }
    }
    & \\ \\

    \theutterance \stepcounter{utterance}  
    & \multicolumn{4}{p{0.6\linewidth}}{
        \cellcolor[rgb]{0.9,0.9,0.9}{
            \makecell[{{p{\linewidth}}}]{
                \texttt{\tiny{[P1$\rangle$GM]}}
                \texttt{STRATEGISCHE ÜBERLEGUNGEN: \{'Der Gegner hat meinen ersten Vorschlag abgelehnt und einen Gegen{-}vorschlag gemacht. Sein Vorschlag enthält A03 (30/594), B61 (97/695), A41 (121/635), C33 (138/313), A04 (262/461), B53 (389/539), B85 (215/4) und A61 (484/203). Gesamtaufwand: 594+695+635+313+461+539+4+203 = 3444, was über dem Maximalaufwand von 3386 liegt! Das ist ein ungültiger Vorschlag. Ich sollte darauf hinweisen, aber taktisch vorgehen. Der Vorschlag enthält einige für mich wertvolle Items (A61, B85, B53) aber auch viele niedrigwertige mit hohem Aufwand. Ich sollte versuchen, die wertvollen Items zu behalten und die schlechten auszutauschen.'\}} \\
\\ 
\texttt{ARGUMENT: \{'Dein Vorschlag überschreitet den Maximalaufwand von 3386 {-} er liegt bei 3444. Das ist nicht zulässig. Zusätzlich enthält er viele Gegenstände mit schlechtem Wert{-}Aufwand{-}Verhältnis. Ich bin bereit zu verhandeln, aber wir müssen eine gültige Lösung finden.'\}} \\
\\ 
\texttt{ABLEHNUNG: \{'A03', 'B61', 'A41', 'C33', 'A04', 'B53', 'B85', 'A61'\}} \\
\\ 
\texttt{VORSCHLAG: \{'C04', 'A61', 'A45', 'B21', 'B85', 'C49', 'A89'\}} \\
            }
        }
    }
    & & \\ \\

    \theutterance \stepcounter{utterance}  
    & & \multicolumn{4}{p{0.6\linewidth}}{
        \cellcolor[rgb]{0.9,0.9,0.9}{
            \makecell[{{p{\linewidth}}}]{
                \texttt{\tiny{[GM$\rangle$P2]}}
                \texttt{ARGUMENT: \{'Dein Vorschlag überschreitet den Maximalaufwand von 3386 {-} er liegt bei 3444. Das ist nicht zulässig. Zusätzlich enthält er viele Gegenstände mit schlechtem Wert{-}Aufwand{-}Verhältnis. Ich bin bereit zu verhandeln, aber wir müssen eine gültige Lösung finden.'\}} \\
\\ 
\texttt{ABLEHNUNG: \{'A03', 'B61', 'A41', 'C33', 'A04', 'B53', 'B85', 'A61'\}} \\
\\ 
\texttt{VORSCHLAG: \{'C04', 'A61', 'A45', 'B21', 'B85', 'C49', 'A89'\}} \\
            }
        }
    }
    & \\ \\

    \theutterance \stepcounter{utterance}  
    & & & \multicolumn{4}{p{0.6\linewidth}}{
        \cellcolor[rgb]{0.9,0.9,0.9}{
            \makecell[{{p{\linewidth}}}]{
                \texttt{\tiny{[GM$\langle$P2]}}
                \texttt{STRATEGISCHE ÜBERLEGUNGEN: \{'Der Gegner hat einen Vorschlag mit C04, A61, A45, B21, B85, C49, A89 gemacht. Aufwand: 313 + 203 + 399 + 437 + 4 + 544 + 720 = 2620. Mein Wert: 30 + 138 + 121 + 262 + 461 + 215 + 65 = 1292. Dieser ist schlechter als mein vorheriger Vorschlag (1810) und enthält A89, das für mich nur 65 Wert hat. Ich sollte auf meinen früheren Vorschlag zurückkommen, aber mit einer kleinen Anpassung, um ihn attraktiver zu machen. Mein optimaler Vorschlag wäre einer, der A03 (668), B61 (508), A41 (500) enthält, aber ich muss den Aufwand im Auge behalten.'\}} \\
\\ 
\texttt{ARGUMENT: \{'Dieser Vorschlag ist für mich weniger vorteilhaft als mein vorheriger. A89 hat für mich einen geringen Wert. Ich schlage vor, wir konzentrieren uns auf Gegenstände mit gutem Wert{-}Aufwand{-}Verhältnis für beide Seiten.'\}} \\
\\ 
\texttt{ABLEHNUNG: \{'C04', 'A61', 'A45', 'B21', 'B85', 'C49', 'A89'\}} \\
\\ 
\texttt{VORSCHLAG: \{'A03', 'B61', 'A41', 'C33', 'A04', 'B53', 'B85'\}} \\
            }
        }
    }
     \\ \\

    \theutterance \stepcounter{utterance}  
    & & \multicolumn{4}{p{0.6\linewidth}}{
        \cellcolor[rgb]{0.9,0.9,0.9}{
            \makecell[{{p{\linewidth}}}]{
                \texttt{\tiny{[P1$\langle$GM]}}
                \texttt{ARGUMENT: \{'Dieser Vorschlag ist für mich weniger vorteilhaft als mein vorheriger. A89 hat für mich einen geringen Wert. Ich schlage vor, wir konzentrieren uns auf Gegenstände mit gutem Wert{-}Aufwand{-}Verhältnis für beide Seiten.'\}} \\
\\ 
\texttt{ABLEHNUNG: \{'C04', 'A61', 'A45', 'B21', 'B85', 'C49', 'A89'\}} \\
\\ 
\texttt{VORSCHLAG: \{'A03', 'B61', 'A41', 'C33', 'A04', 'B53', 'B85'\}} \\
            }
        }
    }
    & \\ \\

    \theutterance \stepcounter{utterance}  
    & \multicolumn{4}{p{0.6\linewidth}}{
        \cellcolor[rgb]{0.9,0.9,0.9}{
            \makecell[{{p{\linewidth}}}]{
                \texttt{\tiny{[P1$\rangle$GM]}}
                \texttt{STRATEGISCHE ÜBERLEGUNGEN: \{'Der Gegner schlägt erneut einen Vorschlag vor, der viele Gegenstände mit schlechtem Wert{-}Aufwand{-}Verhältnis enthält: A03 (30/594), B61 (97/695), A41 (121/635), C33 (138/313), A04 (262/461), B53 (389/539), B85 (215/4). Gesamtaufwand: 594+695+635+313+461+539+4 = 3241, was unter dem Limit liegt. Dieser Vorschlag ist für mich sehr schlecht {-} viele niedrigwertige Items mit hohem Aufwand. Ich sollte klar ablehnen und einen besseren Vorschlag machen, der meine hochwertigen Items enthält.'\}} \\
\\ 
\texttt{ARGUMENT: \{'Dieser Vorschlag enthält überwiegend Gegenstände mit sehr schlechtem Wert{-}Aufwand{-}Verhältnis. Für mich wäre das keine akzeptable Lösung. Ich schlage einen ausgewogeneren Deal vor, der für beide Seiten vorteilhafter ist.'\}} \\
\\ 
\texttt{ABLEHNUNG: \{'A03', 'B61', 'A41', 'C33', 'A04', 'B53', 'B85'\}} \\
\\ 
\texttt{VORSCHLAG: \{'C04', 'A61', 'A45', 'B21', 'C49', 'B85', 'A48'\}} \\
            }
        }
    }
    & & \\ \\

    \theutterance \stepcounter{utterance}  
    & & \multicolumn{4}{p{0.6\linewidth}}{
        \cellcolor[rgb]{0.9,0.9,0.9}{
            \makecell[{{p{\linewidth}}}]{
                \texttt{\tiny{[GM$\rangle$P2]}}
                \texttt{ARGUMENT: \{'Dieser Vorschlag enthält überwiegend Gegenstände mit sehr schlechtem Wert{-}Aufwand{-}Verhältnis. Für mich wäre das keine akzeptable Lösung. Ich schlage einen ausgewogeneren Deal vor, der für beide Seiten vorteilhafter ist.'\}} \\
\\ 
\texttt{ABLEHNUNG: \{'A03', 'B61', 'A41', 'C33', 'A04', 'B53', 'B85'\}} \\
\\ 
\texttt{VORSCHLAG: \{'C04', 'A61', 'A45', 'B21', 'C49', 'B85', 'A48'\}} \\
            }
        }
    }
    & \\ \\

    \theutterance \stepcounter{utterance}  
    & & & \multicolumn{4}{p{0.6\linewidth}}{
        \cellcolor[rgb]{0.9,0.9,0.9}{
            \makecell[{{p{\linewidth}}}]{
                \texttt{\tiny{[GM$\langle$P2]}}
                \texttt{STRATEGISCHE ÜBERLEGUNGEN: \{'Der Gegner schlägt erneut einen ähnlichen Vorschlag vor wie zuvor: C04, A61, A45, B21, C49, B85, A48. Aufwand: 369+203+399+437+544+4+175 = 2131. Mein Wert: 30+138+121+262+215+461+583 = 1810. Das ist der gleiche Wert wie sein erster Vorschlag. Ich muss verstehen, warum er auf diesen Gegenständen besteht {-} sie haben möglicherweise für ihn hohen Wert. Ich könnte versuchen, einige meiner wertvollen Gegenstände hinzuzufügen, während ich einige seiner weniger wertvollen entferne, um den Deal für mich zu verbessern.'\}} \\
\\ 
\texttt{ARGUMENT: \{'Ich sehe, du bevorzugst bestimmte Gegenstände stark. Lass uns einen Kompromiss finden: Ich behalte einige deiner bevorzugten Items, aber füge einige hinzu, die für mich wichtig sind. So können wir beide profitieren.'\}} \\
\\ 
\texttt{ABLEHNUNG: \{'C04', 'A61', 'A45', 'B21', 'C49', 'B85', 'A48'\}} \\
\\ 
\texttt{VORSCHLAG: \{'A03', 'B61', 'A41', 'C33', 'A04', 'B85', 'A48', 'C49'\}} \\
            }
        }
    }
     \\ \\

    \theutterance \stepcounter{utterance}  
    & & \multicolumn{4}{p{0.6\linewidth}}{
        \cellcolor[rgb]{0.9,0.9,0.9}{
            \makecell[{{p{\linewidth}}}]{
                \texttt{\tiny{[P1$\langle$GM]}}
                \texttt{ARGUMENT: \{'Ich sehe, du bevorzugst bestimmte Gegenstände stark. Lass uns einen Kompromiss finden: Ich behalte einige deiner bevorzugten Items, aber füge einige hinzu, die für mich wichtig sind. So können wir beide profitieren.'\}} \\
\\ 
\texttt{ABLEHNUNG: \{'C04', 'A61', 'A45', 'B21', 'C49', 'B85', 'A48'\}} \\
\\ 
\texttt{VORSCHLAG: \{'A03', 'B61', 'A41', 'C33', 'A04', 'B85', 'A48', 'C49'\}} \\
            }
        }
    }
    & \\ \\

    \theutterance \stepcounter{utterance}  
    & \multicolumn{4}{p{0.6\linewidth}}{
        \cellcolor[rgb]{0.9,0.9,0.9}{
            \makecell[{{p{\linewidth}}}]{
                \texttt{\tiny{[P1$\rangle$GM]}}
                \texttt{STRATEGISCHE ÜBERLEGUNGEN: \{'Der Gegner schlägt erneut einen sehr schlechten Deal vor mit vielen niedrigwertigen Items: A03 (30/594), B61 (97/695), A41 (121/635), C33 (138/313), A04 (262/461), B85 (215/4), A48 (65/175), C49 (461/544). Gesamtaufwand: 594+695+635+313+461+4+175+544 = 3421, was immer noch über 3386 liegt! Der Vorschlag ist ungültig und enthält für mich nur einen wirklich wertvollen Gegenstand (C49). Ich muss klar ablehnen und einen besseren Vorschlag machen, der meine Top{-}Items enthält.'\}} \\
\\ 
\texttt{ARGUMENT: \{'Dein Vorschlag überschreitet erneut den Maximalaufwand (3421 > 3386) und enthält zu viele Gegenstände mit geringem Wert. Ich schlage einen fokussierteren Deal vor, der wirklich wertvolle Items für beide Seiten berücksichtigt.'\}} \\
\\ 
\texttt{ABLEHNUNG: \{'A03', 'B61', 'A41', 'C33', 'A04', 'B85', 'A48', 'C49'\}} \\
\\ 
\texttt{VORSCHLAG: \{'C04', 'A61', 'A45', 'B21', 'B85', 'C64', 'A48'\}} \\
            }
        }
    }
    & & \\ \\

    \theutterance \stepcounter{utterance}  
    & & \multicolumn{4}{p{0.6\linewidth}}{
        \cellcolor[rgb]{0.9,0.9,0.9}{
            \makecell[{{p{\linewidth}}}]{
                \texttt{\tiny{[GM$\rangle$P2]}}
                \texttt{ARGUMENT: \{'Dein Vorschlag überschreitet erneut den Maximalaufwand (3421 > 3386) und enthält zu viele Gegenstände mit geringem Wert. Ich schlage einen fokussierteren Deal vor, der wirklich wertvolle Items für beide Seiten berücksichtigt.'\}} \\
\\ 
\texttt{ABLEHNUNG: \{'A03', 'B61', 'A41', 'C33', 'A04', 'B85', 'A48', 'C49'\}} \\
\\ 
\texttt{VORSCHLAG: \{'C04', 'A61', 'A45', 'B21', 'B85', 'C64', 'A48'\}} \\
            }
        }
    }
    & \\ \\

    \theutterance \stepcounter{utterance}  
    & & & \multicolumn{4}{p{0.6\linewidth}}{
        \cellcolor[rgb]{0.9,0.9,0.9}{
            \makecell[{{p{\linewidth}}}]{
                \texttt{\tiny{[GM$\langle$P2]}}
                \texttt{STRATEGISCHE ÜBERLEGUNGEN: \{'Der Gegner schlägt C04, A61, A45, B21, B85, C64, A48 vor. Aufwand: 369+203+399+437+4+685+175 = 2272. Mein Wert: 30+138+121+262+461+97+583 = 1692. Das ist schlechter als vorherige Vorschläge (1692 vs 1810). C64 hat für mich nur 97 Wert. Ich sollte meinen Fokus auf Gegenstände mit hohem Wert{-}Aufwand{-}Verhältnis legen: A03 (668/594=1.125), B61 (508/695=0.731), A41 (500/635=0.787), C33 (484/313=1.546), A04 (400/461=0.868), B53 (389/539=0.722), A48 (583/175=3.331). A48 hat das beste Verhältnis für mich. Ich sollte einen Vorschlag mit A48 und anderen hochwertigen Items machen.'\}} \\
\\ 
\texttt{ARGUMENT: \{'Ich bemerke, dass unsere Vorschläge sich kaum überschneiden. Vielleicht sollten wir offener über unsere Prioritäten sprechen, ohne konkrete Werte preiszugeben. Welche Gegenstände sind dir besonders wichtig?'\}} \\
\\ 
\texttt{ABLEHNUNG: \{'C04', 'A61', 'A45', 'B21', 'B85', 'C64', 'A48'\}} \\
\\ 
\texttt{VORSCHLAG: \{'A03', 'A48', 'C33', 'B61', 'B85', 'A45'\}} \\
            }
        }
    }
     \\ \\

    \theutterance \stepcounter{utterance}  
    & & \multicolumn{4}{p{0.6\linewidth}}{
        \cellcolor[rgb]{0.9,0.9,0.9}{
            \makecell[{{p{\linewidth}}}]{
                \texttt{\tiny{[P1$\langle$GM]}}
                \texttt{ARGUMENT: \{'Ich bemerke, dass unsere Vorschläge sich kaum überschneiden. Vielleicht sollten wir offener über unsere Prioritäten sprechen, ohne konkrete Werte preiszugeben. Welche Gegenstände sind dir besonders wichtig?'\}} \\
\\ 
\texttt{ABLEHNUNG: \{'C04', 'A61', 'A45', 'B21', 'B85', 'C64', 'A48'\}} \\
\\ 
\texttt{VORSCHLAG: \{'A03', 'A48', 'C33', 'B61', 'B85', 'A45'\}} \\
            }
        }
    }
    & \\ \\

    \theutterance \stepcounter{utterance}  
    & \multicolumn{4}{p{0.6\linewidth}}{
        \cellcolor[rgb]{0.9,0.9,0.9}{
            \makecell[{{p{\linewidth}}}]{
                \texttt{\tiny{[P1$\rangle$GM]}}
                \texttt{STRATEGISCHE ÜBERLEGUNGEN: \{'Der Gegner fragt nach Prioritäten, was ich nicht direkt preisgeben darf. Sein neuer Vorschlag ist A03 (30/594), A48 (65/175), C33 (138/313), B61 (97/695), B85 (215/4), A45 (500/399). Gesamtaufwand: 594+175+313+695+4+399 = 2180. Dieser Vorschlag enthält nur einen für mich wertvollen Gegenstand (A45) und viele sehr schlechte Items. Der Aufwand ist sehr niedrig (2180), also gibt es viel Spielraum für Verbesserungen. Ich sollte versuchen, mehr meiner hochwertigen Gegenstände hinzuzufügen.'\}} \\
\\ 
\texttt{ARGUMENT: \{'Ich verstehe deinen Wunsch nach Klarheit über Prioritäten. Für mich sind Gegenstände wichtig, die ein gutes Verhältnis von Wert zu Aufwand bieten. Dein aktueller Vorschlag lässt viel ungenutzten Aufwand und ich würde gerne einige zusätzliche wertvolle Items einbeziehen.'\}} \\
\\ 
\texttt{ABLEHNUNG: \{'A03', 'A48', 'C33', 'B61', 'B85', 'A45'\}} \\
\\ 
\texttt{VORSCHLAG: \{'C04', 'A61', 'A45', 'B21', 'B85', 'C49', 'A89'\}} \\
            }
        }
    }
    & & \\ \\

    \theutterance \stepcounter{utterance}  
    & & \multicolumn{4}{p{0.6\linewidth}}{
        \cellcolor[rgb]{0.9,0.9,0.9}{
            \makecell[{{p{\linewidth}}}]{
                \texttt{\tiny{[GM$\rangle$P2]}}
                \texttt{ARGUMENT: \{'Ich verstehe deinen Wunsch nach Klarheit über Prioritäten. Für mich sind Gegenstände wichtig, die ein gutes Verhältnis von Wert zu Aufwand bieten. Dein aktueller Vorschlag lässt viel ungenutzten Aufwand und ich würde gerne einige zusätzliche wertvolle Items einbeziehen.'\}} \\
\\ 
\texttt{ABLEHNUNG: \{'A03', 'A48', 'C33', 'B61', 'B85', 'A45'\}} \\
\\ 
\texttt{VORSCHLAG: \{'C04', 'A61', 'A45', 'B21', 'B85', 'C49', 'A89'\}} \\
            }
        }
    }
    & \\ \\

    \theutterance \stepcounter{utterance}  
    & & & \multicolumn{4}{p{0.6\linewidth}}{
        \cellcolor[rgb]{0.9,0.9,0.9}{
            \makecell[{{p{\linewidth}}}]{
                \texttt{\tiny{[GM$\langle$P2]}}
                \texttt{STRATEGISCHE ÜBERLEGUNGEN: \{'Der Gegner wiederholt den Vorschlag mit C04, A61, A45, B21, B85, C49, A89 (Aufwand: 369+203+399+437+4+544+720=2676, mein Wert: 30+138+121+262+461+215+65=1292). Dieser ist für mich sehr schlecht, da A89 nur 65 Wert bei 720 Aufwand hat. Ich sollte klar kommunizieren, dass A89 für mich unattraktiv ist, und stattdessen auf Gegenstände mit besserem Wert{-}Aufwand{-}Verhältnis pochen. Mein Fokus liegt auf A48, A03, C33, B61, A41. Vielleicht kann ich einen Vorschlag mit diesen Items machen, plus einigen seiner bevorzugten Items.'\}} \\
\\ 
\texttt{ARGUMENT: \{'A89 hat für mich ein sehr schlechtes Wert{-}Aufwand{-}Verhältnis und wäre in einem Deal für mich nicht akzeptabel. Ich wäre eher bereit, auf andere Gegenstände zu verzichten, wenn wir dafür A48, A03 oder C33 einbeziehen können, die für mich viel wertvoller sind.'\}} \\
\\ 
\texttt{ABLEHNUNG: \{'C04', 'A61', 'A45', 'B21', 'B85', 'C49', 'A89'\}} \\
\\ 
\texttt{VORSCHLAG: \{'A03', 'A48', 'C33', 'B61', 'A41', 'B85', 'A45'\}} \\
            }
        }
    }
     \\ \\

    \theutterance \stepcounter{utterance}  
    & & \multicolumn{4}{p{0.6\linewidth}}{
        \cellcolor[rgb]{0.9,0.9,0.9}{
            \makecell[{{p{\linewidth}}}]{
                \texttt{\tiny{[P1$\langle$GM]}}
                \texttt{ARGUMENT: \{'A89 hat für mich ein sehr schlechtes Wert{-}Aufwand{-}Verhältnis und wäre in einem Deal für mich nicht akzeptabel. Ich wäre eher bereit, auf andere Gegenstände zu verzichten, wenn wir dafür A48, A03 oder C33 einbeziehen können, die für mich viel wertvoller sind.'\}} \\
\\ 
\texttt{ABLEHNUNG: \{'C04', 'A61', 'A45', 'B21', 'B85', 'C49', 'A89'\}} \\
\\ 
\texttt{VORSCHLAG: \{'A03', 'A48', 'C33', 'B61', 'A41', 'B85', 'A45'\}} \\
            }
        }
    }
    & \\ \\

    \theutterance \stepcounter{utterance}  
    & \multicolumn{4}{p{0.6\linewidth}}{
        \cellcolor[rgb]{0.9,0.9,0.9}{
            \makecell[{{p{\linewidth}}}]{
                \texttt{\tiny{[P1$\rangle$GM]}}
                \texttt{STRATEGISCHE ÜBERLEGUNGEN: \{"Der Gegner zeigt klare Präferenzen für A48, A03 und C33 {-} für mich sind das sehr niedrigwertige Items (65, 30, 138), aber er scheint sie zu bevorzugen. Sein Vorschlag A03 (594), A48 (175), C33 (313), B61 (695), A41 (635), B85 (4), A45 (399) hat Gesamtaufwand 594+175+313+695+635+4+399 = 2815. Das ist viel niedriger als das Maximum, also kann ich versuchen, mehr meiner wertvollen Items hinzuzufügen. Ich sollte seine Präferenzen für die niedrigwertigen Items nutzen, um meine hochwertigen Items durchzubekommen."\}} \\
\\ 
\texttt{ARGUMENT: \{"Ich verstehe deine Präferenzen für A48, A03 und C33. Ich bin bereit, diese in den Deal aufzunehmen, möchte aber auch einige meiner wichtigen Gegenstände behalten. Lass uns den Vorschlag erweitern, um den verfügbaren Aufwand besser zu nutzen."\}} \\
\\ 
\texttt{ABLEHNUNG: \{'A03', 'A48', 'C33', 'B61', 'A41', 'B85', 'A45'\}} \\
\\ 
\texttt{VORSCHLAG: \{'C04', 'A61', 'A45', 'B21', 'B85', 'C49', 'A48', 'A03', 'C33'\}} \\
            }
        }
    }
    & & \\ \\

    \theutterance \stepcounter{utterance}  
    & & \multicolumn{4}{p{0.6\linewidth}}{
        \cellcolor[rgb]{0.9,0.9,0.9}{
            \makecell[{{p{\linewidth}}}]{
                \texttt{\tiny{[GM$\rangle$P2]}}
                \texttt{ARGUMENT: \{"Ich verstehe deine Präferenzen für A48, A03 und C33. Ich bin bereit, diese in den Deal aufzunehmen, möchte aber auch einige meiner wichtigen Gegenstände behalten. Lass uns den Vorschlag erweitern, um den verfügbaren Aufwand besser zu nutzen."\}} \\
\\ 
\texttt{ABLEHNUNG: \{'A03', 'A48', 'C33', 'B61', 'A41', 'B85', 'A45'\}} \\
\\ 
\texttt{VORSCHLAG: \{'C04', 'A61', 'A45', 'B21', 'B85', 'C49', 'A48', 'A03', 'C33'\}} \\
            }
        }
    }
    & \\ \\

    \theutterance \stepcounter{utterance}  
    & & & \multicolumn{4}{p{0.6\linewidth}}{
        \cellcolor[rgb]{0.9,0.9,0.9}{
            \makecell[{{p{\linewidth}}}]{
                \texttt{\tiny{[GM$\langle$P2]}}
                \texttt{STRATEGISCHE ÜBERLEGUNGEN: \{'Der Gegner schlägt C04, A61, A45, B21, B85, C49, A48, A03, C33 vor. Aufwand: 369+203+399+437+4+544+175+594+313 = 3042. Mein Wert: 30+138+121+262+461+215+583+668+484 = 2962. Das ist ein sehr guter Wert für mich! Der Aufwand von 3042 ist unter dem Maximum von 3386. Dieser Vorschlag enthält alle meine Top{-}Gegenstände außer B61 und A41. Ich sollte diesen Deal akzeptieren, da er mir fast 3000 Wert bringt und wahrscheinlich für den Gegner auch akzeptabel ist, da er viele seiner bevorzugten Items enthält.'\}} \\
\\ 
\texttt{ARGUMENT: \{'Dieser Vorschlag scheint einen guten Kompromiss zu finden {-} er enthält die für mich wichtigen Gegenstände A48, A03 und C33, während auch deine bevorzugten Items berücksichtigt werden. Der Aufwand bleibt im zulässigen Bereich.'\}} \\
\\ 
\texttt{ZUSTIMMUNG: \{'C04', 'A61', 'A45', 'B21', 'B85', 'C49', 'A48', 'A03', 'C33'\}} \\
            }
        }
    }
     \\ \\

    \theutterance \stepcounter{utterance}  
    & & & \multicolumn{2}{p{0.3\linewidth}}{
        \cellcolor[rgb]{0.9,0.9,0.9}{
            \makecell[{{p{\linewidth}}}]{
                \texttt{\tiny{[GM$|$GM]}}
                \texttt{game successful} \\
            }
        }
    }
    & & \\ \\

    \theutterance \stepcounter{utterance}  
    & & & \multicolumn{2}{p{0.3\linewidth}}{
        \cellcolor[rgb]{0.9,0.9,0.9}{
            \makecell[{{p{\linewidth}}}]{
                \texttt{\tiny{[GM$|$GM]}}
                \texttt{end game} \\
            }
        }
    }
    & & \\ \\

\end{supertabular}
}

\end{document}
